\cotefile{162-5}
\scanpage{36}

\deleted{pleinement fidèle (resp} une équivalence). Même énoncé si on met dans le
coup un \deleted{Anneau \unsure{loc}} systême local d'anneaux sur \(X'\).
E\unsure{\deleted{h}n|\deleted{n}h}fin, \(f\) induit \deleted{une}
\inserted{\deleted{foncteur pl fid (resp} une équivalence)}
\deleted{isomorphisme} entre la catégorie des coefficients locaux sur \(C\) et celle
des coéff. locaux sur \(C'\) \deleted{(?? peut-être cela doit il}
\inserted{\deleted{cette propriété doit}}\deleted{ être misedans les
hypothèses préliminaires sur \(f\) ?)} Principe de démonstration: on commence
par prouver ce dernier résultat, en notant que \deleted{dans le ca} si un topos est
loc\deleted{q}alement hom. trivial, il est loc\unsure{\deleted{a}.|a} connexe et loc. simplement connexe,
d'où une bonne théorie du \(\pi_o\) et du \(\pi_1\) (qui sont ici discrets), et on est
ra\(\unsure{\grave{m}|\overset{\inserted{n}}{m}}\)ené \deleted{pour} à un cas particulier du critère d'homotopisme de Artin et
Mazur, savoir un critère \inserted{cohomologique} pour qu'un homomorhisme de groupes \(G\longrightarrow H\) (ici \deleted{l}des
groupes \(\pi_1\) de \(X,X'\)) soit un isomorphisme: il doit induire des isomorphis\deleted{m}mes
sur les \(H^o\) et \(H^1\) (y inclus dans le cas non commutatif...). On prouve
la pleine fidélité en se ramenant par la suite spectrale aux cas \deleted{\unsure{t|l|e}}des \(\mathrm{Ext}^i\)
de coéfficients locaux. Par la suite spectrale encore, cela résultera du
fait suivant: si \(X^{\unsure{\deleted{\prime}|\prime}}\) est
\unsure{\deleted{h}lo\deleted{m}calement|\deleted{l}ho\deleted{c}malement} hom. trivial, alors la catégorie des
faisceaux abéliens loc. constants est stable par \(\underline{\mathrm{Ext}}^i\), et le foncteur
\(M\longmapsto M_X\) de\deleted{s faisceaux constants} (Ab) dans \(X_{ab}\) commute aux dits \(\mathrm{Ext}^i\).
En fait, \(X\) et \(X'\) étant loc homp\(\mid\)triviaux, les conditions suivantes sur \(f\)
seront équivalentes:

a) \(f\) est\unsure{b} un homotopisme, i.e. induit pour tout systême local (pas
néc commutatif) sur \(X'\)un isomorphisme sur les \(H^i\)\deleted{pour i \unsure{n|m}}.

b) \deleted{Même chose ave} \(f\) induit une équivalence \(D^b_{lc}(X')\longrightarrow D^b_{lc}(X)\).

\(a'\)) \(f\) induit une équivalence entre la catégorie des systêmes locaux
abéliens sur \(X'\) et \(X\), et des isomorphismes sur les \(H^i\) correspondants (donc
on ne prend ici que des coéfficients commutatifs).

\margin{?} J'ignore si on peut dans a) se borner aux sys\deleted{\unsure{y}}têmes locaux commutatifs,
\deleted{i\unsure{.}e\unsure{.} dans}
L'équivalence entre a) et b) fournit une première justification ou
motivation pour définir des types d'homotopie via la catégorie \(D^b_{lc}(X)\),

\margin{35}
\margin{\unsure{r}}
